\section*{Chapter \ref{ch:g12:recursion} --- Recursion}

\begin{solution}{exo:g12:recursion:1}
\smallskip\noindent\hspace*{2em}\code{factorial(3) called}

\noindent\hspace*{4em}\code{factorial(2) called}

\noindent\hspace*{6em}\code{factorial(1) called}

\noindent\hspace*{8em}\code{factorial(0) called}

\noindent\hspace*{8em}\code{factorial(0) returns 1}

\noindent\hspace*{6em}\code{factorial(1) returns 1}

\noindent\hspace*{4em}\code{factorial(2) returns 2}

\noindent\hspace*{2em}\code{factorial(3) returns 6}
\end{solution}

\begin{solution}{exo:g12:recursion:2}
$2^5 - 1 = 31$ moves. The first $15$ move the four smaller disks to B; the $16$th moves disk $5$, the largest, from A
to C; the last $15$ bring the four back on top of it.
\end{solution}

\begin{solution}{exo:g12:recursion:3}
\code{reverse("loop")} is \code{reverse("oop") + "l"}, which is \code{reverse("op") + "o" + "l"}, and so on:
\code{"p" + "o" + "o" + "l"}, that is \code{"pool"}.
\end{solution}

\begin{solution}{exo:g12:recursion:4}
Seven: one call for each length from $6$ down to $0$, the call on the empty list being the \omterm{def:g12:recursion:base-case}{base case} at the top.
(Below them sits the \omterm{def:g12:recursion:call-stack}{frame} of the code that made the first call.)
\end{solution}

\begin{solution}{exo:g12:recursion:5}
\smallskip\noindent\hspace*{2em}\code{def count_digits(n):}

\noindent\hspace*{4em}\code{if n < 10:}

\noindent\hspace*{6em}\code{return 1}

\noindent\hspace*{4em}\code{return 1 + count_digits(n // 10)}

\smallskip\noindent \code{n // 10} drops the last digit. \code{count_digits(2026)} returns $4$; \code{count_digits(0)}
returns $1$, since $0$ is written with one digit.
\end{solution}

\begin{solution}{exo:g12:recursion:6}
There is no \omterm{def:g12:recursion:base-case}{base case}: \code{total(3)} calls \code{total(2)}, \code{total(1)}, \code{total(0)}, \code{total(-1)}, and
so on, until Python stops it with \code{RecursionError}. Correction: \code{if n == 0: return 0} before the \omterm{def:g12:recursion:base-case}{recursive
case}.
\end{solution}

\begin{solution}{exo:g12:recursion:7}
It calls \code{factorial(-2)}, \code{factorial(-3)}, \dots: the \omterm{def:g12:recursion:base-case}{base case} $0$ is never reached, and the run ends with
\code{RecursionError}, a message that does not name the real mistake. Correction, at the top of the function:

\smallskip\noindent\hspace*{2em}\code{if n < 0:}

\noindent\hspace*{4em}\code{raise ValueError("factorial of a negative number")}
\end{solution}

\begin{solution}{exo:g12:recursion:8}
\smallskip\noindent\hspace*{2em}\code{def flatten(nested):}

\noindent\hspace*{4em}\code{if not isinstance(nested, list):}

\noindent\hspace*{6em}\code{return [nested]}

\noindent\hspace*{4em}\code{out = []}

\noindent\hspace*{4em}\code{for item in nested:}

\noindent\hspace*{6em}\code{out += flatten(item)}

\noindent\hspace*{4em}\code{return out}

\smallskip\noindent \code{flatten([1, [2, [3, 4]], [], 5])} returns \code{[1, 2, 3, 4, 5]}.
\end{solution}

\begin{solution}{exo:g12:recursion:9}
Let $P(n)$ be ``\code{power(x, n)} returns $x^n$''. $P(0)$: the function returns $1 = x^0$. If $P(n - 1)$ holds, then
\code{power(x, n)} returns $x \times x^{n-1} = x^n$, so $P(n)$ holds. By induction, $P(n)$ holds for every $n \ge 0$.
It stops: the variant $n$ is a whole number at least $0$, decreased by $1$ at each call, so after $n$ calls it is
$0$, the \omterm{def:g12:recursion:base-case}{base case}.
\end{solution}

\begin{solution}{exo:g12:recursion:10}
\code{1 2 1 3 1 2 1}, one per line. \code{show(n)} prints $p(n) = 2p(n - 1) + 1$ numbers, with $p(0) = 0$: the same
relation as the towers of Hanoi, so $p(n) = 2^n - 1$. (It prints exactly the disks moved by \code{hanoi}, in order.)
\end{solution}

\begin{solution}{exo:g12:recursion:11}
\smallskip\noindent\hspace*{2em}\code{def is_palindrome(s):}

\noindent\hspace*{4em}\code{if len(s) <= 1:}

\noindent\hspace*{6em}\code{return True}

\noindent\hspace*{4em}\code{return s[0] == s[-1] and is_palindrome(s[1:-1])}

\smallskip\noindent The variant is the length, decreased by $2$ at each call; it ends at $0$ or $1$.
\end{solution}

\begin{solution}{exo:g12:recursion:12}
Let $P(n)$: ``every solution for $n$ disks makes at least $2^n - 1$ moves''. $P(1)$: at least one move, $2^1 - 1$.
Assume $P(n - 1)$ and take a solution for $n$ disks. The first time the largest disk moves, from A to some peg, the
other $n - 1$ disks are all on the third peg: it cannot sit on any of them, and it leaves A only when nothing is on
it. Getting them there from A is a solution for $n - 1$ disks: at least $2^{n-1} - 1$ moves, by $P(n - 1)$. The
largest disk moves at least once. After its last move, it lies on C and the $n - 1$ others must come onto it from
wherever they are, all on one peg again, since the largest disk could only have moved with all of them elsewhere: at
least $2^{n-1} - 1$ more moves. In all at least $2(2^{n-1} - 1) + 1 = 2^n - 1$. By induction, $P(n)$ holds for every
$n \ge 1$.
\end{solution}

\begin{solution}{exo:g12:recursion:13}
\code{fib(0)} and \code{fib(1)} make one call each: $c(0) = c(1) = 1$, and $2\,\mathrm{fib}(1) - 1 = 2\,\mathrm{fib}(2)
- 1 = 1$. For $n \ge 2$, \code{fib(n)} makes itself one call plus all the calls of \code{fib(n - 1)} and
\code{fib(n - 2)}: $c(n) = 1 + c(n - 1) + c(n - 2)$. If the formula holds for $n - 1$ and $n - 2$:
$c(n) = 1 + 2\,\mathrm{fib}(n) - 1 + 2\,\mathrm{fib}(n - 1) - 1 = 2\,\mathrm{fib}(n + 1) - 1$. By induction (on two
previous values), it holds for all $n$. Check: $\mathrm{fib}(6) = 8$, $c(5) = 15$; $\mathrm{fib}(26) = 121\,393$,
$c(25) = 242\,785$, as printed.
\end{solution}

\begin{solution}{exo:g12:recursion:14}
\smallskip\noindent\hspace*{2em}\code{def subsets(t):}

\noindent\hspace*{4em}\code{if t == []:}

\noindent\hspace*{6em}\code{return [[]]}

\noindent\hspace*{4em}\code{rest = subsets(t[1:])}

\noindent\hspace*{4em}\code{return rest + [[t[0]] + s for s in rest]}

\smallskip\noindent The subsets of \code{t} are those without \code{t[0]}, the subsets of the rest, and those with
it. If the rest, of $n - 1$ elements, has $2^{n-1}$ subsets, \code{t} has $2 \times 2^{n-1} = 2^n$; the empty list has
$1 = 2^0$. By induction, $2^n$ subsets for $n$ elements.
\end{solution}

\begin{solution}{exo:g12:recursion:15}
\smallskip\noindent\hspace*{2em}\code{def binary(n):}

\noindent\hspace*{4em}\code{if n < 2:}

\noindent\hspace*{6em}\code{return str(n)}

\noindent\hspace*{4em}\texttt{return binary(n // 2) + str(n \% 2)}

\smallskip\noindent Proof, by induction on $n$, assuming the result for every smaller number. For $n = 0$ and $n = 1$,
the writing is the digit itself. For $n \ge 2$, $n = 2q + r$ with $q = \lfloor n/2 \rfloor \ge 1$ and $r = n \bmod 2$. Since
$q < n$, \code{binary(q)} is the binary writing of $q$; appending the digit $r$ multiplies the value by $2$ and adds
$r$: the result is $2q + r = n$, with no leading zero since $q \ge 1$. \code{binary(37)} returns \code{100101}.
\end{solution}

\begin{solution}{pb:g12:recursion:1}
\textbf{1.} The \omterm{def:g12:recursion:base-case}{base case} is $n = 0$: nothing to move. The calls: $n - 1$ disks from \code{start} to \code{spare}
(using \code{goal} as spare), then $n - 1$ disks from \code{spare} to \code{goal} (using \code{start}).

\textbf{2.} When the largest disk moves, the $n - 1$ others are on the spare peg, so its start and goal pegs are
empty; every other move belongs to a smaller puzzle in which disks larger than those moved are all at the bottom of
their pegs, and can carry any of them.

\textbf{3.} Disk $1$: A $\to$ B; disk $2$: A $\to$ C; disk $1$: B $\to$ C.

\textbf{4.} $n + 1$ \omterm{def:g12:recursion:call-stack}{frames} of \code{hanoi}, from $n$ down to $0$: $65$ for $64$ disks, far below the limit of $1000$.

\textbf{5.} The stack is shallow, but the list of moves would hold $2^{64} - 1$ entries, more than any memory can
hold, and producing them would take the time of Part IV.

\textbf{6.} $m(0) = 0$ and $m(n) = 2\,m(n - 1) + 1$.

\textbf{7.} $1$, $3$, $7$, $15$, $31$: as printed.

\textbf{8.} $2^{n-1}$ times: the smallest disk moves at every other move, starting with the first.

\textbf{9.} Once.

\textbf{10.} $m(0) = 0 = 2^0 - 1$. If $m(n - 1) = 2^{n-1} - 1$, then $m(n) = 2(2^{n-1} - 1) + 1 = 2^n - 1$.

\textbf{11.} Exercise \ref{exo:g12:recursion:12}.

\textbf{12.} It makes $2^n - 1$ moves, and no solution makes fewer: it is optimal.

\textbf{13.} That the moves are $2^n - 1$ for $n = 1$ to $10$, with an \code{assert}. Ten cases checked do not prove
the formula for $64$, or for any $n$ not tried; the induction does.

\textbf{14.} Python's integers grow as needed (chapter \ref{ch:g10:machine-integers}); the toy machine's registers
hold one byte.

\textbf{15.} $18\,446\,744\,073\,709\,551\,615$.

\textbf{16.} $365.25 \times 24 \times 3600$: a year of $365$ days and a quarter, the average over leap years.

\textbf{17.} About $584\,542$ million years: $585$ billion, rounded.

\textbf{18.} A billion times fewer: about $585$ years.

\textbf{19.} $80$ years are $2\,524\,608\,000$ seconds. $2^{31} - 1 = 2\,147\,483\,647$ fits, $2^{32} - 1 =
4\,294\,967\,295$ does not: $31$ disks.

\textbf{20.} \textbf{The $64$ disks need $2^{64} - 1 = 18\,446\,744\,073\,709\,551\,615$ moves}, the fewest possible;
\textbf{at one move per second, about $585$ billion years}.
\end{solution}
